\documentclass[troisieme]{fiche}
\metadonnees{5}{Le géant}{Bloc A — Raconter}

\begin{document}
\entete

\objectifs{%
\textbf{Langue} : bilan des temps du bloc ; connecteurs du récit ; thème.\par
\textbf{Texte} : Oscar Wilde, \og The Selfish Giant \fg (1888).\par
\textbf{Durée} : 1 h — exercices 20 min, texte et questions 25 min, version 15 min.}

%% ===================================================================
\exercice[12 min]{Bilan des temps}

\rappel{\textbf{Rappel.} Présent simple pour l'habitude, \emph{be + -ing} pour ce qui se
passe maintenant, prétérit pour ce qui est fini, present perfect pour ce qui compte encore.}

\consigne{Mettez le verbe entre parenthèses au temps qui convient.}

\begin{anglais}
\begin{enumerate}[label=\arabic*.,itemsep=2.5mm,leftmargin=*]
  \item Every afternoon the children \emph{(go)} \trou{went} and played in the garden.
  \item ``How happy we \emph{(be)} \trou{are} here!'' they cried.
  \item The Giant \emph{(stay)} \trou{had stayed} with his friend for seven years.
  \item Look! The children \emph{(climb)} \trou{are climbing} over the wall again.
  \item Spring \emph{(forget)} \trou{has forgotten} this garden, said the Snow.
  \item The birds \emph{(not sing)} \trou{did not sing} in the garden that year.
  \item ``My own garden \emph{(be)} \trou{is} my own garden,'' said the Giant.
  \item He \emph{(build)} \trou{built} a high wall all round it.
  \item The children \emph{(not play)} \trou{have not played} there since the wall went up.
  \item \emph{(you / ever / see)} \trou{Have you ever seen} such a beautiful garden?
\end{enumerate}
\end{anglais}

\note{Prétérit : 1, 6, 8. Présent : 2, 7. \emph{Be + -ing} : 4. Present perfect : 5, 9, 10.
Pluperfect : 3.}

\note[Le partage]{Dans un récit au passé, les paroles gardent le temps où elles ont été
prononcées : d'où le présent des phrases 2 et 7, dites à voix haute. Phrase 3 : le géant est
resté sept ans \emph{avant} de revenir, donc un cran plus loin dans le passé.}

%% ===================================================================
\exercice[8 min]{Thème}
\consigne{Traduisez en anglais.}

\begin{quote}\itshape
Chaque après-midi, les enfants jouaient dans le jardin du géant. Un jour, le géant revint.
Il n'avait pas vu son jardin depuis sept ans. Quand il aperçut les enfants, il cria : \og Que
faites-vous ici ? \fg{} et ils s'enfuirent. Depuis ce jour-là, le jardin est resté vide.
\end{quote}

\lignes{7}

\corr{\textbf{Proposition.} Every afternoon the children played in the Giant's garden. One
day the Giant came back. He had not seen his garden for seven years. When he saw the children,
he cried: ``What are you doing here?'' and they ran away. Since that day the garden has stayed
empty.}

\note[Quatre points]{\og jouaient \fg{} : habitude passée, prétérit simple suffit
(\emph{played}), ou \emph{used to play}. \og depuis sept ans \fg{} après un passé :
\emph{for seven years} avec un pluperfect. \og Que faites-vous ? \fg{} : action en cours,
\emph{what are you doing?} et non \emph{what do you do?} \og Depuis ce jour-là \fg{} + fait
qui dure : present perfect, \emph{has stayed}.}

\note[Le génitif]{\og le jardin du géant \fg{} : \emph{the Giant's garden}. L'anglais met le
possesseur devant, avec \emph{'s}, quand il s'agit d'un être animé.}

%% ===================================================================
\newpage
\section{Texte}

\begin{texte}
Every afternoon, as they were coming from school, the children used to go and play in the
Giant's garden.

It was a large lovely garden, with soft green grass. Here and there over the grass stood
beautiful flowers like stars, and there were twelve peach-trees that in the spring-time broke
out into delicate \gl[a blossom]{blossoms}{une fleur (d'arbre)} of pink and pearl, and in the
autumn \gl[to bear, bore]{bore}{porter, donner} rich fruit. The birds sat on the trees and
sang so sweetly that the children used to stop their games in order to listen to them. ``How
happy we are here!'' they cried to each other.

One day the Giant came back. He had been to visit his friend the Cornish
\gl[an ogre]{ogre}{un ogre}, and had stayed with him for seven years. After the seven years
were over he had said all that he had to say, for his conversation was limited, and he
determined to return to his own castle. When he arrived he saw the children playing in the
garden.

``What are you doing here?'' he cried in a very \gl[gruff]{gruff}{bourru} voice, and the
children ran away.

``My own garden is my own garden,'' said the Giant; ``any one can understand that, and I will
allow nobody to play in it but myself.'' So he built a high wall all round it, and put up a
\gl[a notice-board]{notice-board}{un écriteau}.

\begin{center}\sffamily\small
TRESPASSERS\\ WILL BE\\ PROSECUTED
\end{center}

He was a very selfish Giant.
\end{texte}
\source{Oscar Wilde, \og The Selfish Giant \fg, \emph{The Happy Prince and Other Tales},
Londres, David Nutt, 1888.}

\partiel[15 min]{Questions}
\consigne{Répondez aux deux premières questions en français, aux deux dernières en anglais.}

\begin{enumerate}[label=\arabic*.,itemsep=2.5mm,leftmargin=*]
  \item Décrivez le jardin tel qu'il est au début du conte.
    \lignes{2}
    \corr{Grand et beau, avec une herbe douce et verte, des fleurs éparses comme des étoiles,
    douze pêchers qui fleurissent au printemps et donnent des fruits à l'automne, et des
    oiseaux qui chantent si bien que les enfants s'arrêtent de jouer pour les écouter.}
  \item Pourquoi le géant était-il absent, et pourquoi revient-il ?
    \lignes{2}
    \corr{Il était allé voir son ami l'ogre de Cornouailles et y était resté sept ans. Au
    bout de sept ans il avait dit tout ce qu'il avait à dire — sa conversation était limitée
    — et il décida de rentrer chez lui.}
  \item \en{What does the Giant do when he finds the children in his garden?}
    \lignes{2}
    \corr{He shouts at them in a gruff voice and they run away. Then he builds a high wall
    all round the garden and puts up a notice-board saying that trespassers will be
    prosecuted.}
  \item \en{The last sentence is very short: ``He was a very selfish Giant.'' Who says it,
    and why is it placed there?}
    \lignes{3}
    \corr{The narrator says it, after letting the Giant speak for himself. The judgement
    comes only once the reader has seen the wall and the notice-board, so it sounds like a
    conclusion rather than an opinion.}
\end{enumerate}

\note[\emph{Used to} dans le texte]{\emph{The children used to go and play}, \emph{the
children used to stop their games} : l'habitude passée. Noter \emph{go and play} :
l'anglais coordonne là où le français subordonne (\og aller jouer \fg).}

%% ===================================================================
\section{Version}
\consigne{Traduisez le passage en français. Il suit immédiatement le texte.}

\begin{version}
The poor children had now nowhere to play. They tried to play on the road, but the road was
very \gl[dusty]{dusty}{poussiéreux} and full of hard stones, and they did not like it. They
used to \gl[to wander]{wander}{errer} round the high wall when their lessons were over, and
talk about the beautiful garden inside. ``How happy we were there,'' they said to each other.
\end{version}
\source{\emph{Ibid.}}

\lignes{10}

\corr{\textbf{Proposition de traduction.} Les pauvres enfants n'avaient désormais nulle part
où jouer. Ils essayèrent de jouer sur la route, mais la route était très poussiéreuse et
pleine de pierres dures, et cela ne leur plut pas. Ils prirent l'habitude de rôder autour du
grand mur, leurs leçons finies, et de parler du beau jardin qui se trouvait derrière.
\og Comme nous étions heureux là-bas \fg, se disaient-ils.}

\note[Choix de traduction 1]{\emph{had nowhere to play} : \emph{nowhere} + infinitif se
déplie en français, \og n'avaient nulle part où jouer \fg. Une seule négation en anglais,
deux en français (\og ne\ldots nulle part \fg).}

\note[Choix de traduction 2]{\emph{when their lessons were over} : \emph{to be over} =
\og être fini \fg. \og Leurs leçons finies \fg{} rend le tour plus légèrement que
\og quand leurs leçons étaient finies \fg.}

\note[Choix de traduction 3]{\emph{How happy we were there} : exclamative en \emph{how} +
adjectif, sans inversion. Le français dit \og comme nous étions heureux \fg, et surtout pas
\og comment \fg.}

%% ===================================================================
\partiel{Lexique à mémoriser}
\begin{multicols}{2}\small
\begin{itemize}[label=--,itemsep=0.8mm,leftmargin=*]
  \item \textbf{lovely} — charmant, ravissant
  \item \textbf{grass} — l'herbe
  \item \textbf{a blossom} — une fleur d'arbre
  \item \textbf{sweetly} — doucement
  \item \textbf{a castle} — un château
  \item \textbf{to allow} — permettre
  \item \textbf{a wall} — un mur
  \item \textbf{selfish} — égoïste \textit{(fiche 2 de seconde)}
  \item \textbf{nowhere} — nulle part
  \item \textbf{dusty} — poussiéreux
  \item \textbf{to wander} — errer, rôder
  \item \textbf{to be over} — être fini
\end{itemize}
\end{multicols}

\end{document}
